\chapter{Hai Mặt Của Cuộc Đời (The Two Sides of Life)}
\begin{center}
  \textit{\color{truyengold}``Không có bóng tối, làm sao ta biết ánh sáng quý giá đến nhường nào.''}\\[4pt]
  \textit{(Without darkness, how would we know how precious light truly is.)}\\[4pt]
  \small\color{truyenblue}--- Lão Tử ---
\end{center}
\ngancach

\section{Vua Solomon Và Chiếc Nhẫn}
\begin{truyen}{Vua Solomon Và Chiếc Nhẫn}{Truyện dân gian Do Thái cổ}
\chuhoa{V}{ua} Solomon muốn thử thách vị cố vấn khôn ngoan nhất của mình: ``Hãy tìm cho ta một câu nói có thể làm người buồn vui lên và người vui buồn xuống.''\\[4pt]
\textit{(King Solomon challenged his wisest advisor: `Find me a saying that can cheer a sad man and humble a happy one.')}

Vị cố vấn đi khắp vương quốc, tìm hỏi thợ vàng bạc và học giả.\\[4pt]
\textit{(The advisor traveled throughout the kingdom, consulting goldsmiths and scholars.)}

Cuối cùng ông trở về với một chiếc nhẫn khắc bốn chữ: ``Rồi cũng qua thôi.''\\[4pt]
\textit{(Finally he returned with a ring engraved with four words: `This too shall pass.')}

Vua Solomon nhìn chiếc nhẫn và mỉm cười: khi buồn --- nhớ rằng nỗi buồn rồi sẽ qua; khi vui quá --- nhớ rằng niềm vui cũng rồi sẽ qua; cả hai giúp ta giữ được sự bình thản.\\[4pt]
\textit{(King Solomon looked at the ring and smiled: when sad --- remember the sadness will pass; when too joyful --- remember the joy will also pass; both remind us to maintain equanimity.)}

\end{truyen}
\begin{baihoc}
  Sự bình thản không có nghĩa là không cảm xúc --- mà là hiểu rằng mọi cảm xúc đều tạm thời và mọi hoàn cảnh đều sẽ thay đổi.\\[4pt]
  \textit{(Equanimity does not mean absence of emotion --- it means understanding that all emotions are temporary and all circumstances will change.)}
\end{baihoc}

\section{Đồng Hồ Cát Và Nghịch Lý Thời Gian}
\begin{truyen}{Đồng Hồ Cát Và Nghịch Lý Thời Gian}{Triết lý phương Tây}
\chuhoa{M}{ột} chiếc đồng hồ cát chỉ hoạt động được vì có hai mặt: mặt đầy và mặt vơi; mặt trên và mặt dưới.\\[4pt]
\textit{(An hourglass only works because it has two sides: the full and the empty; the upper and the lower.)}

Nếu ngăn cát chảy xuống --- đồng hồ ngừng. Nếu không có ngăn dưới rỗng để nhận cát --- đồng hồ cũng ngừng.\\[4pt]
\textit{(If you stop the sand from flowing down --- the clock stops. If there is no empty lower chamber to receive the sand --- the clock also stops.)}

Cuộc đời cũng vậy: cho đi mới được nhận lại; mất đi mới có chỗ cho cái mới; đau khổ mới trân trọng được hạnh phúc; thất bại mới hiểu giá trị của thành công.\\[4pt]
\textit{(Life is the same: give to receive; lose to make room for the new; suffer to appreciate happiness; fail to understand the value of success.)}

\end{truyen}
\begin{baihoc}
  Đừng sợ giai đoạn ``vơi'' --- đó là điều kiện cần thiết để giai đoạn ``đầy'' tiếp theo có thể xảy ra.\\[4pt]
  \textit{(Do not fear the `empty' phase --- it is the necessary condition for the next `full' phase to occur.)}
\end{baihoc}

\section{Cô Gái Mù Và Ánh Mặt Trời}
\begin{truyen}{Cô Gái Mù Và Ánh Mặt Trời}{Câu chuyện về Nhận Thức}
\chuhoa{M}{ột} cô gái sinh ra đã mù lần đầu được phẫu thuật phục hồi thị giác. Khi băng mắt được tháo, điều đầu tiên cô thấy là ánh mặt trời buổi sáng xuyên qua cửa sổ.\\[4pt]
\textit{(A girl born blind underwent surgery to restore her sight. When the bandages were removed, the first thing she saw was morning sunlight through the window.)}

Cô òa khóc nức nở --- không phải vì đau, mà vì quá đẹp.\\[4pt]
\textit{(She burst into tears --- not from pain, but because it was too beautiful.)}

Người thân xung quanh hỏi: ``Sao con khóc?'' Cô trả lời: ``Vì trước đây con không biết là thứ mình không có đẹp đến thế nào.''\\[4pt]
\textit{(Family around her asked: `Why are you crying?' She replied: `Because I did not know how beautiful the thing I was missing was.')}

Bài học: chúng ta đang sở hữu hàng ngàn điều quý giá mà ta không nhìn thấy --- vì ta chưa từng mất chúng.\\[4pt]
\textit{(The lesson: we possess thousands of precious things we cannot see --- because we have never lost them.)}

\end{truyen}
\begin{baihoc}
  Lòng biết ơn không đến từ việc có nhiều --- mà từ việc nhìn thấy những gì mình đang có.\\[4pt]
  \textit{(Gratitude does not come from having much --- but from seeing what you already have.)}
\end{baihoc}

\section{Hạt Muối Và Ly Nước --- Hồ Nước}
\begin{truyen}{Hạt Muối Và Ly Nước --- Hồ Nước}{Truyện dạy về Quan Điểm Sống}
\chuhoa{M}{ột} vị thầy tu hòa một nhúm muối vào ly nước nhỏ và bảo học trò uống. Học trò nhăn mặt: ``Đắng quá, thầy ơi.''\\[4pt]
\textit{(A monk dissolved a handful of salt into a small glass of water and told his student to drink it. The student grimaced: `It is too bitter, Master.')}

Vị thầy mang cùng lượng muối đó ra hồ lớn và hòa vào. ``Bây giờ hãy uống nước hồ.'' Học trò uống và lắc đầu: ``Trong vắt, không vị gì cả.''\\[4pt]
\textit{(The monk took the same amount of salt to a large lake and dissolved it. `Now drink the lake water.' The student drank and shook his head: `Crystal clear, no taste at all.')}

Vị thầy nói: ``Nỗi đau trong cuộc sống giống như muối --- lượng muối không thay đổi, nhưng cái thay đổi là dung tích của bình chứa. Khi bạn đau khổ, hãy mở rộng lòng ra, không phải thu mình lại.''\\[4pt]
\textit{(The monk said: `The pain in life is like salt --- the amount of salt does not change, but what changes is the capacity of the container. When you suffer, expand your heart, do not contract it.')}

\end{truyen}
\begin{baihoc}
  Nỗi đau không thay đổi theo ý ta --- nhưng tâm hồn ta có thể lớn hơn nỗi đau.\\[4pt]
  \textit{(Pain does not change at our will --- but our soul can grow larger than the pain.)}
\end{baihoc}
